\documentclass[11pt]{article}

\usepackage[a4paper,margin=28mm]{geometry}
\usepackage[T1]{fontenc}
\usepackage{lmodern}
\usepackage{microtype}
\usepackage{amsmath,amssymb,amsthm,mathtools}
\usepackage{booktabs}
\usepackage{array}
\usepackage{enumitem}
\usepackage{xcolor}
\usepackage[hidelinks]{hyperref}
\setlength{\emergencystretch}{2em}

\newtheorem{theorem}{Theorem}[section]
\newtheorem{proposition}[theorem]{Proposition}
\newtheorem{corollary}[theorem]{Corollary}
\newtheorem{lemma}[theorem]{Lemma}
\theoremstyle{definition}
\newtheorem{definition}[theorem]{Definition}
\theoremstyle{remark}
\newtheorem{remark}[theorem]{Remark}

\newcommand{\F}{\mathbb{F}}
\newcommand{\cC}{\mathcal{C}}
\newcommand{\cH}{\mathcal{H}}
\newcommand{\cR}{\mathcal{R}}
\newcommand{\cS}{\mathcal{S}}
\newcommand{\cT}{\mathcal{T}}
\newcommand{\supp}{\operatorname{supp}}
\newcommand{\wt}{\operatorname{wt}}
\newcommand{\Tr}{\operatorname{Tr}}
\newcommand{\rank}{\operatorname{rank}}
\newcommand{\Span}{\operatorname{span}}

\title{Complete repair hypergraphs under concatenation\\
\large A twisted-cubic--axis family}
\author{Tavis Rudd}
\date{July 2026}

\begin{document}
\maketitle

\begin{abstract}
For a linear code $C$ and a coordinate $x$, let $\cH_x^{(r)}(C)$ be the
complete hypergraph of all linear recovery sets of size at most $r$ at $x$.
Its matching number $\nu_x$ is disjoint
availability, while its transversal number $\tau_x$ is the least number of
helper failures that block every radius-$r$ repair.  These invariants retain
overlap information that locality and availability alone discard.

For $q=3^h$, we study the $[2q+1,4,q-1]_q$ code generated by the $q$
affine points of a twisted cubic together with the $q+1$ points of its
characteristic-three axis.  Axis coordinates have exact locality two and
cubic coordinates exact locality three.  We determine the axis invariants
in terms of the affine cap number, prove
$\left(\nu_x,\tau_x\right)=\left((q-1)/2,q-2\right)$ at every cubic
coordinate, and obtain $\tau_x>\nu_x$ for every
coordinate when $q\geq9$.  At $q=9$ the three coordinate types have exact rows
\[
 (\nu,\tau)=(4,7),\qquad(6,12),\qquad(7,13).
\]
Restoring the cubic point at infinity gives a more symmetric $[2q+2,4,q]_q$
variant whose full inclusion-minimal repair port stabilizes at radius four,
with uniform rows
\[
 \left(\frac{q-1}{2},q-1\right)\quad\hbox{on the cubic},
 \qquad
 \left(\frac{5q-3}{6},2q-3\right)\quad\hbox{on the axis}.
\]

Our second result is an exact concatenation theorem.  If $I$ is an inner code and
$r+1<2d(I^\perp)$, an outer functional-dual distance of at least $r+2$
forces every concatenated dual word of weight at most $r+1$ into one inner
block.  Consequently the entire hypergraph $\cH_x^{(r)}$, rather than a
chosen repair subfamily, is copied block by block.  A trace-duality argument
turns ordinary extension-field dual distance into this functional gate.
Combining the $q=9$ seed with a self-dual TVZ outer family gives unbounded
$\F_9$-linear codes of exact rate $2/19$ whose relative distance is eventually
greater than every fixed $c<39/190$ (in particular, greater than $1/5$),
while preserving the exact repair rows blockwise.  The completed seed gives
a rate-$1/10$ family with eventual relative distance above every
$c<351/1600$ and exact radius-four rows $(4,8)$ and $(7,15)$.  These are
bounded-port statements.  All finite arguments and reductions from the sole
deep input, Stichtenoth's self-dual TVZ theorem, are formalized in Lean.
\end{abstract}

\section{Introduction}

Locality asks how many surviving symbols are sufficient to recover one lost
coordinate of a code.  Availability refines locality by asking for many
pairwise-disjoint recovery sets.  These parameters have become standard in
the theory of locally repairable codes; see, for example,
Gopalan--Huang--Simitci--Yekhanin~\cite{GopalanEtAl2012} and the geometric
constructions of Pamies-Juarez--Hollmann--Oggier~\cite{PamiesEtAl2013} and
Lavrauw--Van de Voorde~\cite{LavrauwVanDeVoorde2019}.

Disjoint availability does not determine robustness against adversarial
helper failures.  The full family of recovery sets may overlap in a way that
admits a small blocking set, or it may remain repairable after every set of
failures much larger than a maximum disjoint family.  The natural object is
therefore the complete bounded repair hypergraph at a coordinate.  We write
$\nu$ for its matching number and $\tau$ for its transversal number.  Then
$\nu$ is disjoint availability, while $\tau-1$ is the maximum number of
arbitrary helper failures that can always be tolerated while retaining at
least one local repair.  Our emphasis is the exact coordinatewise computation
of both invariants for a complete code-derived repair family and their
simultaneous transfer.  The transversal interpretation itself is established:
Pamies-Juarez--Hollmann--Oggier~\cite[Definition~3]{PamiesEtAl2013} introduced
the same minimum hitting-set invariant as local repair tolerance, and
Wang--Zhang~\cite{WangZhang2014} placed it in a general regenerating-set
framework.

Our inner families come from a classical object in finite projective
geometry.  In $\mathrm{PG}(3,q)$, take the affine part of a twisted cubic and,
when $q$ has characteristic three, adjoin the common axis of the osculating
planes.  We study both this punctured system and the system obtained by
restoring the projective cubic point at infinity.  The choice is driven by
its small dependencies: triples on the axis give two-helper repairs, while
pairs of cubic points complete with an axis point to give three-helper
repairs.  The axis and its
stabilizer orbits are classical; a modern account
appears in Davydov--Marcugini--Pambianco~\cite{DavydovEtAl2021}.  What matters
here is not an orbit classification but the complete system of small linear
dependencies through each selected point.

The second ingredient is concatenation.  Concatenation with a locally
repairable inner code is established machinery and yields asymptotically
good LRC families; Liu--Ma--Wu--Xing~\cite{LiuEtAl2022}, in particular,
derive locality and asymptotic bounds this way.  Locality preservation only
asserts that some inner repairs survive.  We isolate a stronger dual-distance
condition that prevents new low-weight dual words from spanning several
inner blocks.  Under this condition, \emph{every} bounded repair support in
the concatenated code is an embedded inner repair support.  Thus matching,
transversal, edge counts, and all other invariants of the complete repair
hypergraph transfer exactly.

The argument therefore has two arcs.  Sections~\ref{sec:repair}--\ref{sec:projective-completion}
construct the repair hypergraph and compute it for the punctured and completed
geometric seeds.  Sections~\ref{sec:transfer}--\ref{sec:asymptotic} prove exact
blockwise transfer, connect its functional hypothesis to ordinary dual
distance, and derive two fixed-alphabet asymptotic families.

\subsection*{Contributions}

The main contributions follow the same two-part structure.

\begin{enumerate}[label=(\roman*),leftmargin=2.1em]
\item We construct a characteristic-three code $C_q$ with parameters
      $[2q+1,4,q-1]_q$ and classify every minimal linear dependency of size
      at most four.
\item We compute exact axis-coordinate matching and transversal numbers in
      terms of the affine cap number (an exact reduction to a standard finite
      additive invariant), prove the exact cubic-coordinate row
      $((q-1)/2,q-2)$, and deduce $\tau_x>\nu_x$ for every
      coordinate when $q\geq9$.  The case $q=9$, used below as the inner
      seed, has exact rows $(4,7)$, $(6,12)$, and $(7,13)$.
\item We restore the projective cubic point at infinity, prove parameters
      $[2q+2,4,q]_q$, and determine the complete inclusion-minimal repair
      port.  Every inclusion-minimal repair has at most four helpers, with
      uniform cubic row $((q-1)/2,q-1)$ and axis row
      $((5q-3)/6,2q-3)$.
\item We prove exact preservation of complete radius-$r$ repair hypergraphs
      under the block-confinement hypothesis
      $r+1<2d(I^\perp)$ and an outer functional-dual distance gate.  As a
      derived corollary, every exact inner locality $s\leq r$ is preserved.
      Explicit $\F_3$ examples show that both numerical thresholds are best
      possible for a theorem uniform over all inner and outer codes.
\item We show that ordinary extension-field dual distance implies the
      functional gate, by representing component functionals through the
      trace pairing.
\item Using Stichtenoth's self-dual TVZ family~\cite{Stichtenoth2005}, we
      derive a fixed-alphabet family over $\F_9$ with exact rate $2/19$,
      eventual relative distance greater than every fixed $c<39/190$, exact
      mixed locality, exact type multiplicities, exact repair rows, and
      helper-failure thresholds.
      The projectively completed seed gives a second, Pareto-incomparable
      point with exact rate $1/10$, every eventual distance constant below
      $351/1600$, equal $10N/10N$ coordinate multiplicities, and exact
      radius-four rows $(4,8)$ and $(7,15)$.
\end{enumerate}

\subsection*{Scope of the contribution}

The underlying geometry, repair tolerance as a hitting number, and
locality-preserving concatenation are prior art.  The contributions studied
here are narrower: exact all-symbol matching and transversal numbers for this
point system, and equality of its \emph{complete} bounded repair hypergraph
under an outer-dual gate.  Section~\ref{sec:prior} gives the detailed
comparison with adjacent work.

\section{Complete bounded repair hypergraphs}
\label{sec:repair}

Let $C\leq\F_q^X$ be a linear code on a finite coordinate set $X$, with
dual $C^\perp$ under the standard dot product.

\begin{definition}\label{def:repair-hypergraph}
Fix $x\in X$ and $r\geq0$.  The \emph{complete radius-$r$ repair
hypergraph} at $x$ records exactly the helper support of every dual relation
involving $x$ and at most $r$ other coordinates:
\[
 \cH_x^{(r)}(C)=
 \left\{R\subseteq X\setminus\{x\}: |R|\leq r,
 \begin{array}{l}
 \text{there is }y\in C^\perp\text{ with }y_x\ne0,\\[-2mm]
 \supp(y)=\{x\}\cup R
 \end{array}\right\}.
\]
Its inclusion-minimal members form the repair clutter
$\cH_{x,\min}^{(r)}(C)$.  We put
\[
 \nu_x^{(r)}(C)=\nu(\cH_x^{(r)}(C)),\qquad
 \tau_x^{(r)}(C)=\tau(\cH_x^{(r)}(C)),
\]
where $\nu$ is maximum matching size and $\tau$ is minimum transversal size.
\end{definition}

The word $y$ supplies the recovery equation
\[
 c_x=-y_x^{-1}\sum_{z\in R}y_zc_z\qquad(c\in C).
\]
Thus Definition~\ref{def:repair-hypergraph} includes every linear repair
whose coefficient support has at most $r$ helpers; it is not a declared or
selected subfamily.

\begin{proposition}\label{prop:basic-hypergraph}
For every finite code-derived repair hypergraph with nonempty edges,
\[
 \nu_x^{(r)}(C)\leq\tau_x^{(r)}(C).
\]
Passing to the inclusion-minimal repair clutter preserves both $\nu$ and
$\tau$.  The dual distance also controls edge uniformity: if
$r+1\leq d(C^\perp)$, every edge has exactly $r$ helpers.
\end{proposition}

\begin{proof}
A transversal meets every edge of a matching, and the matching edges are
disjoint, proving $\nu\leq\tau$.  Every edge contains an inclusion-minimal
edge.  Consequently a set hits all edges exactly when it hits all minimal
edges; replacing each matching edge by a minimal subedge preserves
disjointness.  Finally, an edge $R$ is witnessed by a nonzero dual word of
weight $|R|+1$.  Dual distance gives $r+1\leq |R|+1$, while the definition
gives $|R|\leq r$.
\end{proof}

The operational meaning of $\tau$ is immediate: every failure set of size
at most $\tau-1$ misses some repair edge, whereas a transversal of size
$\tau$ blocks them all.  Matching and transversal can nevertheless be far
apart; equality is not a generic hypergraph law.

\begin{remark}[Operational scope]\label{rem:operational-scope}
The displayed equation gives one direct scalar protocol: each helper
forwards its stored $\F$-symbol, so access and download are both $|R|$ field
symbols.  This is achieved scalar help-by-transfer, not a bandwidth or access
optimum in a subpacketized model~\cite{ShahEtAl2010,GuruswamiWootters2016,YeBarg2016}.
Nor are the raw coefficients invariants: diagonal column rescaling changes
their ratios while preserving every dual support.  The results below transfer
support invariants such as disjoint availability and failure tolerance, not
optimal bandwidth, optimal access, or arithmetic cost.
\end{remark}

\section{The characteristic-three twisted-cubic--axis code}
\label{sec:seed}

We seek a projective system with two complementary kinds of small
dependencies: collinear triples for locality two, and mixed
cubic--axis circuits for locality three.  Characteristic three supplies both
in one classical configuration.

Let $q=3^h$ and write $\F=\F_q$.  In homogeneous coordinates on
$\mathrm{PG}(3,q)$ define
\[
 T_q=\{C(t)=(1:t:t^2:t^3):t\in\F\}
\]
and
\[
 L_q=\{A(u)=(0:1:u:0):u\in\F\}\cup\{A_\infty=(0:0:1:0)\}.
\]
The set $T_q$ is the affine twisted cubic and $L_q$ is the full line
$X_0=X_3=0$.  For $q\equiv0\pmod3$, this is the common axis of the
osculating-plane pencil~\cite[Section~7]{DavydovEtAl2021}.  We use that
geometric name only as provenance: all proofs below start from the displayed
coordinates.

Let $S_q=T_q\sqcup L_q$ and let $G_q$ be the $4\times(2q+1)$ matrix whose
columns are fixed representatives of these points.  Write $C_q$ for the row
code of $G_q$.  Changing a projective representative only rescales one code
coordinate, so the repair supports are unaffected.

\begin{theorem}[Code parameters]\label{thm:parameters}
For every finite field of characteristic three,
\[
 C_q=[2q+1,4,q-1]_q.
\]
\end{theorem}

\begin{proof}
The length is $2q+1$, and, for example,
$C(0),C(1),A(0),A_\infty$ span $\F^4$.  The distance
of a projective-system code is its length minus the largest hyperplane
section.  If a plane does not contain $L_q$, it meets $L_q$ in at most one
point and its equation restricts to a nonzero cubic polynomial on $T_q$, so
it contains at most three points of $T_q$.  If it contains $L_q$, its
equation has the form $aX_0+dX_3=0$.  On $T_q$ this becomes $a+dt^3=0$,
which has at most one solution because Frobenius $t\mapsto t^3$ is
injective.  Hence every plane section has size at most $q+2$.  The plane
$X_3=0$ contains all of $L_q$ and the point $C(0)$, attaining $q+2$.
Therefore $d=(2q+1)-(q+2)=q-1$.
\end{proof}

\subsection{Small circuits and exact locality}

For distinct $s,t,u\in\F$, put
\[
 e_1=s+t+u,\qquad e_2=st+su+tu.
\]
These symmetric sums are precisely the coordinates of the axis point that
completes three cubic points to a dependency.

\begin{theorem}[Small-circuit classification]\label{thm:circuits}
The inclusion-minimal dependent subsets of $S_q$ having size at most four
are exactly:
\begin{enumerate}[label=(\alph*),leftmargin=2em]
\item the triples of distinct points of $L_q$;
\item the sets
\[
 \{C(s),C(t),C(u),(0:e_1:e_2:0)\}
\]
for distinct $s,t,u\in\F$.
\end{enumerate}
The pair $(e_1,e_2)$ is never $(0,0)$ for distinct $s,t,u$, so the axis
completion in (b) is a unique projective point.  Consequently every axis
coordinate has exact locality two and every cubic coordinate has exact
locality three.
\end{theorem}

\begin{proof}
Every three distinct points on $L_q$ are minimally dependent.  Vandermonde
independence excludes all-cubic sets of size at most four because the cubic
point at infinity has been punctured.  Among mixed four-sets, a direct
$4\times4$ determinant is nonzero for two cubic and two axis columns.  For
three cubic columns and one axis column, expansion gives dependence exactly
at $(0:e_1:e_2:0)$, establishing both existence and uniqueness.
If $e_1=e_2=0$, then $s,t,u$ would be the three roots of
$X^3-stu$; injectivity of Frobenius would force them equal.  Thus the four
points form a circuit.  The locality assertions follow from the
dual-support/circuit correspondence and the absence of smaller circuits
through the respective target type.
\end{proof}

\section{Matching and transversal invariants}
\label{sec:invariants}

Let $Z_3(q)$ denote the largest size of a subset of the additive group of
$\F_q$ containing no three distinct elements with sum zero.  Under an
$\F_3$-basis, this is the usual affine cap number in $\mathrm{AG}(h,3)$.
It enters because complements of axis-repair transversals are precisely such
zero-sum-free sets.

\subsection{Axis coordinates}

Fix $y\in L_q$.  The minimal radius-three repair clutter at $y$ splits over
two disjoint helper grounds, one containing only axis helpers and the other
only cubic helpers:
\begin{itemize}[leftmargin=2em]
\item every pair among the $q$ other axis points repairs $y$, giving the
      complete graph $K_q$;
\item a cubic triple repairs $y$ exactly when its completion in
      Theorem~\ref{thm:circuits} equals $y$.
\end{itemize}
Matching and transversal numbers add across this disjoint union.

For $A_\infty$, the cubic triples are exactly the distinct zero-sum triples
of $\F_q$, hence the affine lines of $\mathrm{AG}(h,3)$.  For $A(a)$, the
following projective relabeling moves the target to $A_\infty$:
\[
 t\longmapsto(t+a)^{-1},\qquad 0^{-1}=0,
\]
The convention at the pole $t=-a$ makes this a permutation of $\F_q$.
Under it, the cubic repair component becomes the zero-sum triples avoiding
zero.

\begin{theorem}[Exact axis rows]\label{thm:axis-rows}
At the point $A_\infty$,
\[
 \nu_{A_\infty}=\frac{5q-3}{6},
 \qquad
 \tau_{A_\infty}=2q-1-Z_3(q).
\]
At every finite axis point $A(a)$,
\[
 \nu_{A(a)}=\frac{5q-9}{6},
 \qquad
 \tau_{A(a)}=2q-2-Z_3(q).
\]
These are the invariants of both the minimal repair clutter and the complete
radius-three repair hypergraph.
\end{theorem}

\begin{proof}
The complete graph $K_q$ contributes matching $(q-1)/2$ and transversal
$q-1$.  The affine-line hypergraph contributes matching $q/3$ and
transversal $q-Z_3(q)$: parallel lines give the matching, and complements
exchange transversals with zero-sum-free sets, since a complement contains no
repair triple exactly when the removed set hits every one.  After forbidding
zero, one parallel class loses its line through zero and the largest
zero-sum-free complement is unchanged; the values become $q/3-1$ and
$q-1-Z_3(q)$.  Adding the contributions gives the formulas.
Proposition~\ref{prop:basic-hypergraph} passes the invariants between the
complete hypergraph and its minimal clutter.
\end{proof}

The strict gap does not require a quantitative cap-set bound.  Any matching
in the minimal axis clutter consumes two axis helpers for each pair edge and
three cubic helpers for each cubic edge.  Weighting axis helpers by three and
cubic helpers by two assigns total weight six to either edge type, and yields
$6\nu_y\leq5q$.  The $K_q$ component alone forces
$\tau_y\geq q-1$.  Thus $\nu_y<\tau_y$ for $q\geq9$.

\subsection{Cubic coordinates}

Fix $x\in\F$.  By Theorem~\ref{thm:circuits}, every unordered pair
$\{s,t\}\subseteq\F\setminus\{x\}$ gives one repair
\[
 \{C(s),C(t),A(\kappa_x(s,t))\},
\]
where $\kappa_x(s,t)$ is the unique completion color.  For fixed $s$, the
map $t\mapsto\kappa_x(s,t)$ is injective away from the diagonal.  The repair
hypergraph is therefore a properly edge-colored $K_{q-1}$ in which every
graph edge is augmented by its color vertex on the axis.

\begin{theorem}[Exact cubic rows]\label{thm:cubic-rows}
Every cubic coordinate has exactly $\binom{q-1}{2}$ radius-three repairs and
\[
 \nu_x=\frac{q-1}{2},\qquad
 \tau_x=q-2.
\]
In particular, $\tau_x>\nu_x$ for $q\geq9$.
\end{theorem}

\begin{proof}
The edge count is the number of unordered pairs of cubic helpers.  A matching
uses two distinct cubic helpers per edge, giving $\nu_x\leq(q-1)/2$.  For the
reverse inequality, we relabel the cubic helpers so completion colors are
injectively determined by sums, then choose a perfect matching with distinct
sums.  Put
\[
 u=(s-x)^{-1},\qquad v=(t-x)^{-1}.
\]
As $s,t$ range over the cubic helpers, $u,v$ range over $\F_q^*$.  Direct
simplification in characteristic three shows that the completing axis point
is
\[
 \begin{cases}
 A_\infty,&u+v=0,\\
 A\bigl(-x+(u+v)^{-1}\bigr),&u+v\ne0.
 \end{cases}
\]
This is an injective recoding of the sum $u+v$.  Choose a generator $g$ of
$\F_q^*$ and pair $g^{2i}$ with $g^{2i+1}$ for
$0\leq i<(q-1)/2$.  These pairs partition $\F_q^*$, while their sums
$(1+g)g^{2i}$ are distinct.  For $q=3$ there is only one pair, so the
conclusion is immediate even though its sum is zero.  Mapping back by
$s=x+u^{-1}$ gives a rainbow perfect matching of the $q-1$ cubic helpers and
therefore $\nu_x=(q-1)/2$.

For the transversal upper bound, choose one cubic helper $a\ne x$ and take
all cubic helpers other than $x$ and $a$; this set has size $q-2$ and every
repair uses one of them.  Conversely, suppose a transversal leaves a
nonempty set $U$ of cubic helpers unchosen.  Fix $a\in U$.  For each
$b\in U\setminus\{a\}$, the edge on $\{a,b\}$ forces its completion color
into the transversal, and properness makes these colors distinct.  Counting
the chosen cubic helpers together with these forced colors gives at least
$q-2$ vertices.
\end{proof}

Combining the two coordinate types gives the uniform statement.

\begin{corollary}[All-symbol separation]\label{cor:uniform-gap}
For every $q=3^h\geq9$ and every coordinate $z$ of $C_q$,
\[
 \nu_z^{(3)}(C_q)<\tau_z^{(3)}(C_q).
\]
Moreover, $C_q$ has all-symbol locality three, with exact locality two on
$L_q$ and exact locality three on $T_q$.
\end{corollary}

\subsection{The exact field-nine seed}

For $q=9$, the affine cap number is $Z_3(9)=4$, and the general cubic formula
gives $\nu_x=4$.  We obtain the complete table below; repair counts refer to
inclusion-minimal repairs, while $(\nu,\tau)$ is unchanged for the complete
hypergraph.

\begin{table}[ht]
\centering
\begin{tabular}{@{}lccccl@{}}
\toprule
coordinate type & multiplicity & locality & minimal repairs & $\nu$ & $\tau$ \\
\midrule
cubic $C(x)$       & $9$ & $3$ & $28$      & $4$ & $7$ \\
finite axis $A(a)$ & $9$ & $2$ & $36+8$    & $6$ & $12$ \\
axis $A_\infty$    & $1$ & $2$ & $36+12$   & $7$ & $13$ \\
\bottomrule
\end{tabular}
\caption{Exact rows of the $[19,4,8]_9$ seed.}
\label{tab:q9}
\end{table}

The $36$ pair repairs are the edges of $K_9$.  The cubic components contain
eight zero-sum triples avoiding zero at a finite axis point and twelve
zero-sum triples at $A_\infty$.  Globally, the small-circuit inventory consists of
$\binom{10}{3}=120$ axis triples and $\binom93=84$ completed cubic
four-circuits.

These rows also give the uniform ratio
\[
 7\nu_z\leq4\tau_z.
\]
The constant $7/4$ is attained at every cubic coordinate.

\section{The projectively completed seed}
\label{sec:projective-completion}

Restoring the missing cubic point has two effects: it recovers projective
symmetry within each parameter line, and it makes the complete
inclusion-minimal repair port accessible through the ambient rank-four
bound.

Let
\[
 \overline T_q=T_q\cup\{C_\infty=(0:0:0:1)\},
 \qquad \overline S_q=\overline T_q\sqcup L_q,
\]
and let $\overline C_q$ be the row code of the displayed projective system.
Thus $\overline C_q$ restores the one cubic point punctured from $C_q$; the
axis is unchanged.

\begin{theorem}[Completed parameters and full minimal repair port]
\label{thm:completed-port}
For every $q=3^h$,
\[
 \overline C_q=[2q+2,4,q]_q,
 \qquad d(\overline C_q^\perp)=3.
\]
Cubic coordinates have exact locality three and axis coordinates exact
locality two.  Moreover, every inclusion-minimal repair uses at most four
helpers, so the full minimal repair port is already visible at radius four.
Its exact rows are
\[
 (\nu_x^{(4)},\tau_x^{(4)})=
 \begin{cases}
 ((q-1)/2,q-1),&x\in\overline T_q,\\[2mm]
 ((5q-3)/6,2q-3),&x\in L_q.
 \end{cases}
\]
These values are uniform within each of the two projective parameter lines.
\end{theorem}

\begin{proof}
A plane not containing the axis meets it once and meets the full twisted
cubic in at most three points.  A nonzero plane containing the axis is an
osculating plane and meets the full cubic in exactly one point.  Hence every
section has size at most $q+2$, and the old extremal section still attains
$q+2$.  This gives distance $(2q+2)-(q+2)=q$; the same spanning columns give
dimension four.  Pairwise projective independence gives dual distance at
least three, while any three distinct axis points give equality.

The new boundary four-circuit
\[
 \{C(s),C(t),C_\infty,A(s+t)\}\qquad(s\ne t)
\]
supplies the cubic-infinity repairs.  The shifted-inversion relabeling used
above extends projectively across its pole.  Its ambient linear realization
permutes the selected columns up to nonzero scaling and moves any chosen
cubic, respectively axis, target to the corresponding infinity target.
Thus locality and the repair hypergraph need be calculated at only those two
targets.

For exhaustiveness, take an inclusion-minimal target-plus-helper relation.
Its helper columns are linearly independent: otherwise a helper-only
relation can be subtracted to cancel one helper without cancelling the
target coefficient, because that helper-only relation has target coefficient
zero.  Since the generator has four rows, there are at most
four helpers.  This proves that radius four is the full \emph{minimal} inner
port; it makes no assertion about a later concatenated code's unbounded
port.

For a cubic target, every minimal edge consumes at least two of the $q$
other cubic points, so $\nu\le(q-1)/2$; the radius-three rainbow matching
attains equality.  For an axis target, if an edge uses $c$ cubic and $a$
axis helpers, the short-circuit exclusions imply: if $a=0$ then $c\ge3$; if
$a=1$ then $c\ge2$; and if $a\ge2$ the inequality is immediate.  Hence
$2c+3a\ge6$.  Summing over a matching
and using the $q+1$ cubic and $q$ available axis resources gives
$6\nu\le5q+2$, hence $\nu\le(5q-3)/6$; the radius-three matching attains
equality.

Finally, the complement of a full-port transversal contains no helper set
that spans the target.  It therefore lies in a hyperplane avoiding the
target; conversely every such hyperplane complement is a transversal.
Equivalently,
\[
 \tau_x=(2q+1)-m_x,
\]
where $m_x$ is the largest hyperplane section avoiding $x$.  The maximum is
$q+2$ for a cubic target and $4$ for an axis target, giving $q-1$ and
$2q-3$, respectively.
\end{proof}

At $q=9$, the completed seed is $[20,4,9]_9$.  Its ten cubic coordinates
all have radius-four row $(4,8)$, while its ten axis coordinates all have
row $(7,15)$.  The two old affine axis subtypes merge because the restored
projective symmetry acts transitively on the full axis.

\section{Exact repair transfer under concatenation}
\label{sec:transfer}

Informally, each outer symbol is a vector encoded as one inner block.  The
question is whether a small dual relation can couple several such blocks.
We formulate the answer without choosing coordinates on the outer symbol
space.  Let $I\leq
\F^K$ be an inner code and let $V$ be an $\F$-vector space equipped with a
linear isomorphism $e:V\to I$.  For an outer $\F$-linear code
$O\leq V^J$, the concatenated code $O\circ_e I\leq\F^{J\times K}$ encodes
each outer symbol through $e$.

The \emph{functional dual} of $O$ is
\[
 O^{\perp_{\rm fun}}=
 \left\{(\beta_j)_{j\in J}\in(V^*)^J:
       \sum_{j\in J}\beta_j(u_j)=0\text{ for every }u\in O\right\}.
\]
Its distance is the minimum number of nonzero component functionals in a
nonzero tuple; its support therefore counts the outer blocks touched by the
induced dual relation.  This is intrinsic to $V$ and does not choose field
coordinates.

\begin{theorem}[Complete repair-hypergraph transfer]\label{thm:transfer}
Let $r\geq0$.  Suppose
\[
 r+1<2d(I^\perp)
 \quad\text{and}\quad
 d(O^{\perp_{\rm fun}})\geq r+2.
\]
Then for every block $j\in J$ and inner coordinate $x\in K$,
\[
 \cH_{(j,x)}^{(r)}(O\circ_e I)
 =\{\{j\}\times R:R\in\cH_x^{(r)}(I)\}.
\]
Consequently matching number, transversal number, inclusion-minimal repair
clutter, edge counts, and the full incidence structure are preserved
coordinate by coordinate.
\end{theorem}

The first inequality is the \emph{inner block-confinement gate}: two nonzero
inner-dual blocks would already be too heavy.  The second is the
\emph{outer functional gate}: a relation touching at most $r+1$ blocks must
induce the zero functional on each one.

\begin{proof}
Take $w\in(O\circ_e I)^\perp$ with $\wt(w)\leq r+1$ and write $w_j$ for
its blocks.  Each $w_j$ defines a functional $\beta_j\in V^*$ by pairing
$w_j$ with $e(v)$.  Orthogonality to the concatenated code says
$(\beta_j)_j\in O^{\perp_{\rm fun}}$.  Its functional support has size at
most $\wt(w)\leq r+1$, below the outer functional-dual distance.  Hence
every $\beta_j$ is zero, so each block $w_j$ lies in $I^\perp$.

Every nonzero block now has weight at least $d(I^\perp)$.  If two blocks
were nonzero, then
\[
 \wt(w)\geq2d(I^\perp)>r+1,
\]
a contradiction.  Thus a bounded dual word containing $(j,x)$ is supported
only in block $j$, where its helper support is an inner repair of $x$.  This
proves one inclusion.  Conversely, extending any inner dual word by zero
outside block $j$ gives a dual word of every concatenation, proving the
other inclusion.
\end{proof}

\begin{corollary}[Exact mixed-locality transfer]\label{cor:exact-locality-transfer}
Under the hypotheses of Theorem~\ref{thm:transfer}, let $0\leq s\leq r$.
For every $(j,x)$, the concatenated coordinate has exact locality $s$ if
and only if the inner coordinate $x$ has exact locality $s$.
\end{corollary}

\begin{proof}
The inner gate at radius $s$ follows from $s+1\leq r+1<2d(I^\perp)$, and
the outer gate follows from $d(O^{\perp_{\rm fun}})\geq r+2\geq s+2$.
Apply Theorem~\ref{thm:transfer} at every radius at most $s$.  Existence of
an $s$-repair and nonexistence at every smaller radius are therefore both
preserved.
\end{proof}

The inner gate does not require
$d(I^\perp)=r+1$.  For the full $q=9$ seed, $d(C_9^\perp)=3$ because of
the axis triples, yet radius-three transfer is exact since $4<2\cdot3$.

\subsection{Sharpness of the transfer gates}

\begin{proposition}[Uniform boundary of the transfer gates]
\label{prop:transfer-boundary}
Neither numerical hypothesis in Theorem~\ref{thm:transfer} can be weakened
uniformly, even among nonzero codes over $\F_3$.
\begin{enumerate}[label=(\alph*),leftmargin=2.1em]
\item At radius $r=3$ there are an inner code and a full outer code with
      $r+1=2d(I^\perp)$ for which complete repair-hypergraph equality fails.
\item At radius $r=2$ there are an inner code satisfying
      $r+1<2d(I^\perp)$ and an outer code with functional-dual distance
      exactly $r+1$ for which complete repair-hypergraph equality fails.
\end{enumerate}
Thus the strict inner inequality and the outer threshold $r+2$ are best
possible for the stated uniform theorem.  This does not assert that either
hypothesis is necessary for a fixed concatenation.
\end{proposition}

\begin{proof}
For both examples let $I\leq\F_3^2$ be the repetition code
$I=\{(a,a):a\in\F_3\}$.  Its dual is generated by $(1,-1)$, so
$d(I^\perp)=2$.

For (a), use two outer symbols and take the full outer space.  Its functional
dual is zero, so it satisfies every outer gate.  The concatenated dual word
\[
 (1,-1\mid 1,-1)
\]
has weight four and meets both blocks.  With its first coordinate as target,
the other three support coordinates form a radius-three repair which is not
an embedded repair of one inner block.  Here $4=2d(I^\perp)$.

For (b), use three outer symbols and the single-parity-check outer code
\[
 O=\{(a,b,c)\in\F_3^3:a+b+c=0\}.
\]
Its functional dual consists of constant triples of scalar functionals and
has exact functional-dual distance three.  In the concatenation, the word
that is one in the first coordinate of each block and zero elsewhere is a
dual word of weight three.  Targeting its first coordinate gives a
two-helper repair meeting all three blocks, so it is not an embedded inner
repair.  Meanwhile $3<2d(I^\perp)=4$.

Both examples give strict failure of complete repair-hypergraph equality.
The two-block dual-distance mechanism is the same elementary symmetric-
difference obstruction used for distinct recovery sets, for example by
Kurz--Yaakobi~\cite[Lemma~10(a)]{KurzYaakobi2021}; the point here is the exact
uniform boundary of the complete-hypergraph transfer statement.
\end{proof}

\section{The trace bridge and the field-nine lift}
\label{sec:trace}

Let $L/\F$ be a finite separable extension and let $O\leq L^J$ be an
$L$-linear outer code.  Restricting scalars makes each outer symbol an
$\F$-vector space.  The functional-dual gate in
Theorem~\ref{thm:transfer} is then supplied by the trace pairing.

\begin{theorem}[Trace-duality bridge]\label{thm:trace}
Every $\F$-linear functional $f:L\to\F$ has a unique representation
\[
 f(x)=\Tr_{L/\F}(a_fx).
\]
If $(f_j)_j$ belongs to the functional dual of the restricted-scalar code,
then $(a_{f_j})_j\in O^\perp$.  Moreover
\[
 f_j=0\quad\Longleftrightarrow\quad a_{f_j}=0
\]
at every coordinate.  Hence
\[
 d\bigl((O|_\F)^{\perp_{\rm fun}}\bigr)\geq d(O^\perp).
\]
\end{theorem}

\begin{proof}
Separable trace is a nondegenerate bilinear pairing, giving the unique
coefficient $a_f$ and exact support preservation.  If $u\in O$ and
$c\in L$, functional orthogonality applied to $cu\in O$ gives
\[
 0=\sum_j f_j(cu_j)
  =\Tr_{L/\F}\!\left(c\sum_j a_{f_j}u_j\right)
 \qquad\text{for every }c\in L.
\]
Nondegeneracy forces $\sum_ja_{f_j}u_j=0$.  Thus the coefficient tuple lies
in the ordinary $L$-dual of $O$, with the same support.
\end{proof}

Now take $\F=\F_9$, $L=\F_{9^4}=\F_{6561}$, and use a linear equivalence
between the four-dimensional $\F_9$-space $L$ and the four-dimensional inner
code $C_9$.

\begin{corollary}[Finite extension-field lift]\label{cor:finite-lift}
Let $O\leq L^N$ be an $[N,K,D]_L$ code with $O\ne0$ and
$d(O^\perp)\geq5$.  Its concatenation with $C_9$ is an $\F_9$-linear code
with
\[
 [19N,\ 4K,\ \geq8D]_9.
\]
Its complete radius-three repair hypergraph is an embedded copy of the
corresponding $C_9$ hypergraph in each block.  Consequently, among its $19N$
coordinates, the exact blockwise profile is:
\[
\begin{array}{c|c|c|c}
\text{type}&\text{multiplicity}&\text{locality}&(\nu,\tau)\\ \hline
\text{cubic}&9N&3&(4,7)\\
\text{finite axis}&9N&2&(6,12)\\
\text{infinity axis}&N&2&(7,13).
\end{array}
\]
The guaranteed arbitrary-helper failure thresholds $\tau-1$ are
respectively $6$, $11$, and $12$.
\end{corollary}

\begin{proof}
Restriction of scalars multiplies dimension by four and does not change
symbol support.  A nonzero outer word has at least $D$ nonzero symbols, each
of whose inner encodings has weight at least eight, so the concatenated
distance is at least $8D$.  Theorem~\ref{thm:trace} supplies functional-dual distance at
least five, while $4<2d(C_9^\perp)=6$.  Apply
Theorem~\ref{thm:transfer}, Corollary~\ref{cor:exact-locality-transfer},
and Table~\ref{tab:q9}.
\end{proof}

\begin{corollary}[Completed finite extension-field lift]
\label{cor:completed-finite-lift}
Let $O\leq L^N$ be an $[N,K,D]_L$ code with $O\ne0$ and
$d(O^\perp)\geq6$.  Concatenation with $\overline C_9$ is an $\F_9$-linear
code with
\[
 [20N,\ 4K,\ \geq9D]_9.
\]
Its complete repair hypergraphs through radius four are embedded copies of
the corresponding inner hypergraphs.  Exactly $10N$ coordinates are cubic,
with exact locality three and row $(4,8)$ at radius four; exactly $10N$ are
axis coordinates, with exact locality two and row $(7,15)$.  The guaranteed
helper-failure thresholds are $7$ and $14$.
\end{corollary}

\begin{proof}
The parameter argument is the same, now using the $[20,4,9]_9$ inner seed.
At radius four the inner gate is $5<2d(\overline C_9^\perp)=6$, and the trace
bridge requires ordinary outer dual distance $r+2=6$.  Theorem~\ref{thm:transfer}
therefore gives equality of the complete bounded hypergraphs.  Apply
Theorem~\ref{thm:completed-port}.  As in Theorem~\ref{thm:transfer}, the
conclusion concerns the radius-four port.
\end{proof}

\section{A fixed-alphabet asymptotic family}
\label{sec:asymptotic}

We require one deep existence theorem.
Stichtenoth~\cite[Theorem~1.6(ii)]{Stichtenoth2005} proves that for every
square $Q=\ell^2$ there is an unbounded sequence of self-dual
$[N_j,N_j/2,D_j]_Q$ codes with
\[
 \liminf_j\frac{D_j}{N_j}\geq
 \frac12-\frac1{\ell-1}.
\]
For $Q=6561=81^2$, the lower bound is $39/80$.  Self-duality is especially
useful here: the same $D_j$ supplies both the primal distance needed for the
concatenated distance and the dual distance needed for exact repair
transfer.

\begin{theorem}[Uniform asymptotic repair family]\label{thm:asymptotic}
There is a sequence of $\F_9$-linear codes $\cC_j$ with lengths tending to
infinity such that, for all sufficiently large $j$,
\begin{enumerate}[label=(\alph*),leftmargin=2em]
\item $\cC_j$ has exact rate $2/19$;
\item its relative minimum distance is at least $1/5$;
\item in every inner block, nine cubic coordinates have exact locality
      three, nine finite-axis coordinates exact locality two, and one
      infinity-axis coordinate exact locality two;
\item on those three respective types, the complete radius-three repair
      rows are $(4,7)$, $(6,12)$, and $(7,13)$, and the guaranteed
      helper-failure thresholds are $6$, $11$, and $12$.
\end{enumerate}
In particular $7\nu_x\leq4\tau_x$ at every coordinate.
More sharply, for every fixed real $c<39/190$, one has
\[
 \frac{d(\cC_j)}{\operatorname{length}(\cC_j)}>c
\]
for all sufficiently large $j$ (with the threshold allowed to depend on
$c$).
\end{theorem}

\begin{proof}
Take Stichtenoth's self-dual codes over $L=\F_{6561}$ and concatenate them
with $C_9$ after restriction to $\F_9$.  Two numerical consequences are
needed: dual distance at least five for the transfer gate, and enough primal
distance for the displayed relative-distance bound.  Since $39/80>1/3$,
eventually
\[
 N_j\leq3D_j
 \quad\text{and}\quad N_j\geq15.
\]
The second inequality and the first imply $D_j\geq5$, so self-duality gives
$d(O_j^\perp)=D_j\geq5$ and Corollary~\ref{cor:finite-lift} applies.
Since also $39/80>19/40$, eventually $19N_j\leq40D_j$.
The lifted length and dimension are $19N_j$ and
$4(N_j/2)=2N_j$, giving exact rate $2/19$.  Its distance is at least
$8D_j$, and
\[
\frac{d(\cC_j)}{19N_j}
 \geq\frac{8D_j}{19N_j}
 \geq\frac15.
\]
The exact localities and repair rows transfer blockwise.
For the sharper distance statement, fix $c<39/190$.  Then
$19c/8<39/80$, so eventually $D_j/N_j>19c/8$.  Hence
\[
 \frac{d(\cC_j)}{19N_j}
 \geq\frac{8D_j}{19N_j}>c.
\]
\end{proof}

\begin{theorem}[Completed uniform asymptotic repair family]
\label{thm:completed-asymptotic}
There is a sequence of $\F_9$-linear codes $\overline\cC_j$ of unbounded
length and exact rate $1/10$ such that, eventually, the relative minimum
distance is at least $1/5$.  Each inner block has ten cubic coordinates of
exact locality three and radius-four row $(4,8)$, and ten axis coordinates
of exact locality two and row $(7,15)$.  More sharply, for every fixed real
$c<351/1600$,
\[
 \frac{d(\overline\cC_j)}{\operatorname{length}(\overline\cC_j)}>c
\]
eventually, with the threshold allowed to depend on $c$.
\end{theorem}

\begin{proof}
Use the same self-dual outer family and concatenate with $\overline C_9$.
Here the corresponding goals are dual distance at least six and the
completed distance ratio.  Eventually $N_j\ge18$ and $N_j\le3D_j$, hence
$D_j\ge6$ and
Corollary~\ref{cor:completed-finite-lift} applies.  Length and dimension are
$20N_j$ and $4(N_j/2)=2N_j$, giving rate $1/10$.  The distance is at least
$9D_j$.  Since $39/80>4/9$, eventually $4N_j<9D_j$, which yields the clean
$1/5$ bound.  For fixed $c<351/1600=9(39/80)/20$, one has
$20c/9<39/80$, so eventually
\[
 \frac{d(\overline\cC_j)}{20N_j}
 \geq\frac{9D_j}{20N_j}>c.
\]
The bounded radius-four hypergraph equality transfers the coordinate
profiles.  The conclusion remains the bounded radius-four equality supplied
by Theorem~\ref{thm:transfer}.
\end{proof}

The tradeoff between the two constructions is summarized in
Table~\ref{tab:family-comparison}.  Completion slightly lowers the rate,
raises the available asymptotic distance threshold, merges the axis
subtypes, and transfers one additional repair radius.

\begin{table}[ht]
\centering
\begin{tabular}{@{}lcc@{}}
\toprule
 & punctured seed & completed seed \\
\midrule
inner parameters & $[19,4,8]_9$ & $[20,4,9]_9$ \\
exact asymptotic rate & $2/19$ & $1/10$ \\
proved eventual $c$ & $c<39/190$ & $c<351/1600$ \\
transferred radius & $3$ & $4$ \\
block multiplicities & $9+9+1$ & $10+10$ \\
repair rows & $(4,7),(6,12),(7,13)$ & $(4,8),(7,15)$ \\
\bottomrule
\end{tabular}
\caption{Comparison of the two fixed-alphabet families.  The distance row
lists every eventual strict relative-distance constant proved above.}
\label{tab:family-comparison}
\end{table}

\section{Relation to prior work}
\label{sec:prior}

The locality of codeword symbols and its parameter bounds originate in the
modern LRC literature represented by~\cite{GopalanEtAl2012}.  Multiple
repair alternatives and finite-geometric constructions are explicit in
Pamies-Juarez--Hollmann--Oggier~\cite{PamiesEtAl2013}; exact availability for
large classes of generalized quadrangles is developed in
Lavrauw--Van de Voorde~\cite{LavrauwVanDeVoorde2019}.  In particular,
\cite[Definition~3]{PamiesEtAl2013} defines local repair tolerance by the
same minimum hitting-set formula as our $\tau$, while Wang--Zhang
\cite{WangZhang2014} organize several tolerance notions through regenerating
sets.  Thus neither repair tolerance nor the hypergraph interpretation is
claimed as new here.  The results developed here are the exact
coordinatewise pair $(\nu,\tau)$, its strict separation on every symbol of
this family, and the transfer of the entire bounded repair family.

Gruica--Jany--Ravagnani~\cite{GruicaEtAl2026} give a refined weight
distribution indexed by prescribed subsets of dual-word supports and prove
duality identities for it.  This is close prior art for treating all bounded
dual supports rather than a selected repair group.  It does not state the
matching/transversal computations here or a concatenation theorem that
identifies the complete bounded support family.

For concatenation, Liu--Ma--Wu--Xing
\cite[Theorem~3.1]{LiuEtAl2022} concatenate a locally repairable inner code
with classical outer codes and derive locality and asymptotic parameter
families.  The recent construction of Jin--Fu~\cite{JinFu2026} also uses
concatenation to obtain binary LRCs with disjoint local repair groups and
controlled weight distributions.  These results establish locality and code
parameters, not equality of every bounded repair support.  Our transfer
theorem strengthens the conclusion under an additional outer-dual gate: it
identifies every bounded dual support and therefore transfers the complete
repair hypergraph, not only the existence of an inner recovery set.

The geometry of the twisted cubic and its stabilizer has a substantial
literature.  Davydov--Marcugini--Pambianco~\cite{DavydovEtAl2021} describe
the axis in characteristic three and line orbits under the cubic stabilizer;
Bartoli--Davydov--Marcugini--Pambianco~\cite{BartoliEtAl2020} compute
point--plane incidence data and multiple-covering-code parameters.  Those
incidence counts are adjacent to, but do not by themselves determine,
matching and transversal numbers of the repair hypergraphs in
Sections~\ref{sec:invariants} and~\ref{sec:projective-completion}.  The full
projective cubic and its characteristic-three common axis are therefore not
claimed as new.  The new calculation concerns the exact full-minimal repair
rows of their union, not the underlying geometry or ordinary code parameters.

The multiplicative pairing used in the cubic matching proof belongs to a
classical design-theoretic neighborhood.  Constant-quotient pairings of a
finite-field half-set are studied as one-quotient starters; see, for example,
Dinitz~\cite{Dinitz1984}.  We do not claim this pairing pattern as new.  The
use made here is to combine it with the shifted-inverse completion
identity to compute the exact matching number of the code-derived repair
hypergraph.

In summary, the contribution is the conjunction of an explicit all-symbol
$\tau>\nu$ seed, exact complete-hypergraph transfer, and the resulting
fixed-alphabet parameters.

\section{Formal verification and provenance}
\label{sec:verification}

The project was human-directed and agent-generated.  The named author
selected the problem, directed the research, set the validation standards,
and takes responsibility for the final claims; AI systems generated
substantial portions of the mathematics, code, formal proofs, and prose.
AI is not listed as an author.

The complete finite theorem chain is formalized in Lean 4 using mathlib.
The formal objects are ordinary finite-dimensional linear codes,
dual-word supports, finite hypergraphs, matching and transversal numbers,
finite fields, restriction of scalars, and the algebra trace.  The formal
proofs establish four linked groups of results:
\begin{itemize}[leftmargin=2em]
\item the finite geometry, code parameters, small-circuit classification,
      exact localities, and every matching/transversal row;
\item complete bounded repair-hypergraph transfer, preservation of exact
      locality through the transferred radius, and both sharpness examples;
\item the support-preserving trace bridge, the concrete degree-four field
      extension, and both finite lifted-code profiles;
\item every asymptotic reduction, including the two exact rates, the strict
      distance thresholds $39/190$ and $351/1600$, and the common displayed
      bound $1/5$.
\end{itemize}

No \texttt{sorry}, \texttt{admit}, \texttt{native\_decide},
\texttt{unsafe}, or code-generation trust occurs in this proof chain.
The finite checks use kernel reduction.  Axiom reports for all finite,
algebraic, and analytic-reduction declarations contain only
\texttt{propext}, \texttt{Classical.choice}, and \texttt{Quot.sound}.
The asymptotic headline adds exactly one project axiom,
\path{Imported.stichtenoth_selfDual_TVZ_6561}, whose docstring cites
Stichtenoth's theorem number, specialization, and exact parameter statement.
Thus the asymptotic theorems are ordinary consequences of Stichtenoth's
published result; within the project trust ledger, that exact specialization
is the single quarantined literature axiom.

The trust manifest is
\texttt{lean/RepairCodes/TRUST.md}.  The paper package includes a separate
proof ledger mapping every displayed theorem to Lean declarations and an
adversarial novelty review separating mathematical novelty from
formalization novelty.

\appendix
\section{Statement-adequacy excerpts}
\label{app:adequacy}

This appendix gives typographically normalized transcriptions of the
load-bearing Lean declarations.  Typeclass hypotheses---finite fields,
finite index types, decidable equality, and finite-dimensionality---are
declared immediately above these statements in the source files.

\subsection*{Exact cubic matching number}
\begin{verbatim}
theorem cubicRepair_matchingNumber [CharP F 3] (x : F) :
    matchingNumber (axisTwistedCubicRepairHypergraph (.inl x) 3) =
      (Fintype.card F - 1) / 2
\end{verbatim}
Its lower bound uses the separately kernel-checked declarations
\begin{verbatim}
units_addColor_matchingNumber_lower
twistedCubicTripleAxisIndex_eq_cubicRepairSumColor
\end{verbatim}
for the finite-field matching and shifted-inverse completion identity.

\subsection*{Exact complete-repair transfer}
\begin{verbatim}
theorem repairHypergraph_concatenatedCode_eq_embed
    (I : Submodule F (kappa -> F)) (e : V ~=l[F] I)
    (O : Submodule F (iota -> V)) (r : Nat)
    (hsI : r + 1 < 2 * dualDist I)
    (hdO : HasFunctionalDualDistanceAtLeast O (r + 2))
    (j : iota) (x : kappa) :
    repairHypergraph (concatenatedCode I e O) (j, x) r =
      embedHypergraph (blockEmbedding j) (repairHypergraph I x r)
\end{verbatim}
The ASCII substitutions \texttt{F}, \texttt{kappa}, \texttt{iota}, and
\texttt{\~=l} normalize the Unicode identifiers in the checked source.

The checked corollary
\texttt{hasExactLocalityAt\_concatenatedCode\_iff\_of\_le} states that the
same radius-$r$ gates preserve the predicate \texttt{HasExactLocalityAt} at
every $s\leq r$.

\subsection*{Trace-duality bridge}
\begin{verbatim}
theorem hasFunctionalDualDistanceAtLeast_restrictScalars
    (O : Submodule L (iota -> L)) (d : Nat)
    (hd : d <= dualDist O) :
    HasFunctionalDualDistanceAtLeast (O.restrictScalars K) d
\end{verbatim}

\subsection*{Exact extension lift}
\begin{verbatim}
theorem q9ExtensionLiftCode_parameters
    (hcard : Fintype.card F = 9)
    (hdeg : Module.finrank F L = 4)
    (O : Submodule L (iota -> L)) (hO : O != bottom) :
    Fintype.card (iota x AxisTwistedCubicIndex F) =
        19 * Fintype.card iota /\
    Module.finrank F (q9ExtensionLiftCode hdeg O) =
        4 * Module.finrank L O /\
    8 * minDist O <= minDist (q9ExtensionLiftCode hdeg O)
\end{verbatim}

\subsection*{Concrete asymptotic headline}
\begin{verbatim}
theorem concrete_q9_uniform_repair_family :
    HasQ9UniformRepairFamily finrank_q9OuterField :=
  stichtenoth_q9_uniform_repair_family
    card_q9BaseField card_q9OuterField finrank_q9OuterField
\end{verbatim}
Here \texttt{HasQ9UniformRepairFamily} expands to unbounded length, eventual
$19k=2n$, eventual $n\leq5d$, every strict relative-distance constant below
$39/190$, a disjoint exhaustive coordinate partition with multiplicities
$9N,9N,N$, exact locality three/two, the exact three-way matching/transversal
row match, and thresholds $6,11,12$.  The complete definition is recorded in
\texttt{RepairCodes/Asymptotic.lean}; the external proof ledger points to
the exact source line and axiom report.

\bibliographystyle{plain}
\bibliography{refs}

\end{document}
